\section{The literature by constrained party}
\label{sec:literature}

Section~\ref{sec:findings-tables} reports how the corpus distributes. This section
says what is in it. The order follows the party each region constrains, because
that is the axis on which the regions differ, and each subsection states what the
region contains, the shape its characteristic result takes, and the boundary that
shape imposes.

\subsection{The attacker bounded in time}

The largest region bounds an adversary in running time or circuit size,
\AttackerTimeCirc{} works in all. Its characteristic result is a guarantee that
holds against every efficient algorithm and against no unbounded one.

Computational differential privacy is the definitional core. Mironov et al.\ relax
the differential privacy guarantee from every distinguisher to every efficient one,
producing the indistinguishability and simulation variants that later work takes as
given \cite{MPRV09}. What that relaxation buys is contested. Groce, Katz and
Yerukhimovich show it buys nothing in the single-curator setting, though their
result is narrower than it is usually quoted as being, applying to black-box
constructions and to queries whose output is a tuple of real numbers of constant
dimension \cite{GKY11}. The authors exclude, by name, queries returning
classifiers, graphs or synthetic databases, which is the case a released model
falls into.

Robustness against a bounded attacker is the region's constructive side. Garg, Jha,
Mahloujifar and Mahmoody wrap a learning problem whose classifier is vulnerable to
every perturbation into one whose classifier resists polynomial-time adversaries
and still yields to unbounded ones, using signatures and error-correcting codes
\cite{GJMM20}. Etesami, Mahloujifar and Mahmoody give a polynomial-time algorithm
that finds, on product distributions under Hamming distance, the close points the
concentration-of-measure bound promises, so there the geometric inevitability of
Section~\ref{sec:lit-none} is attained by an efficient attacker \cite{EMM20}.
Cohen et al.\ define a computational form of reconstruction resilience against
polynomial-time attackers and show that, for training maps that are efficiently
computable and well-spread, it is equivalent to one-wayness \cite{CKMMNST25}.

\subsection{The attacker bounded by what it may change}

The next region, \AttackerCapability{} works, bounds what the adversary may alter
and leaves what it may spend open. \AttackerPerturbAll{} fix a perturbation radius
and \AttackerCorruptAll{} a fraction of the training set.

Certified defences are the mature part. Cohen, Rosenfeld and Kolter certify a radius
around an input within which a smoothed classifier's prediction cannot change
\cite{CRK19}, and Steinhardt, Koh and Liang certify a bound on the loss a poisoning
adversary can induce while controlling a fixed fraction of the data \cite{SKL17}.
Both are guarantees against an adversary with unlimited computation and a limited
reach. The evaluation critiques sit in this region too. Athalye et al.\ and
Tram\`er et al.\ break defences inside their own threat models, and the threat
model each defence stated fixes a radius and nothing else \cite{ACW18,TCBM20}.

That combination is the region's defining feature and its boundary. A radius says
what an attack may do and nothing about what it costs, so two defences certified at
the same radius are indistinguishable under it even when one is far more expensive
to attack. Section~\ref{sec:findings-tables} records that this is the most common
form of adversary bound in the attack literature.

\subsection{The learner bounded}

\LearnerBound{} works constrain the party that trains rather than the party that
attacks. \LearnerTime{} bound it in time, \LearnerSamples{} in samples,
\LearnerClass{} in an algorithm class, \LearnerFlops{} in a budget of
floating-point operations, and \LearnerOther{} in parameters, circuit size, memory
or queries.

The algorithm-class results are the ones that transfer. Feldman et al.\ prove that
statistical-query algorithms need superpolynomially many queries to detect a planted
clique \cite{FGRVX17}, and the low-degree framework gives the same shape of bound
against polynomial thresholds \cite{Wei25,BBHLS21}. Because gradient descent and
principal component analysis are statistical-query algorithms, these bounds apply to
the methods learning actually uses. Bubeck, Price and Razenshteyn use exactly this
to separate robust from ordinary learning, and their separation rests on no
cryptographic assumption at all \cite{BPR19}.

The time-bounded results include the one that runs the implication from robustness
to cryptography. Degwekar, Nakkiran and Vaikuntanathan prove that a task where an
efficient robust classifier exists but cannot be learned efficiently implies
one-way functions \cite{DNV19}. The direction matters and is easy to invert. The
theorem licenses no conclusion from the belief that one-way functions exist, and
what it supports is that a learning task with that hardness profile carries
cryptographic structure rather than arising by accident.

Sample complexity is the second unit. Montasser, Hanneke and Srebro show that every
class of finite VC dimension can be learned under adversarial perturbation, though
not by empirical risk minimization and sometimes at a sample cost exponential in
the VC dimension \cite{MHS19}, and later characterize the sample complexity of the
agnostic case \cite{MHS22}. Nothing in either is computational.

Two results bound the learner in a resource the rest of the corpus does not use.
Bubeck and Sellke show that smoothly interpolating $n$ points in $d$ dimensions
needs about $nd$ parameters, an information-theoretic law with no adversary in it
\cite{BS21}. Raz shows that learning parity from a stream needs either quadratic
memory or exponentially many samples \cite{Raz16}, which is the memory bound the
bounded-storage literature builds on \cite{GZ19}.

\subsection{Nobody bounded}
\label{sec:lit-none}

\NumUnstated{} works constrain no party at all. They are not a residue, and two
distinct lines sit inside them.

The first concludes inevitability from geometry. Mahloujifar, Diochnos and Mahmoody
prove that in a concentrated metric probability space any classifier with constant
error is vulnerable to small perturbations \cite{MDM19}, and Gilmer et al.\ exhibit
the effect on a constructed distribution of two concentric spheres
\cite{GMFSRWG18}. The radius in each theorem is its conclusion, not a restriction
on any adversary. These conclude inevitability precisely because no resource enters
them, which is why the frame of this survey does not hold them and why the survey
says so.

The hypothesis those theorems rest on has been measured. Mahloujifar, Zhang,
Mahmoody and Evans estimate the concentration function on image benchmarks and find
it does not explain observed vulnerability \cite{MZME19}. A conditional theorem
whose antecedent is measured false is a limit on the theorem, not a sibling result,
and the corpus records it as one.

The second line is critique of evaluation and of definition. Das et al.\ show that
classifiers which never query a model beat published membership inference attacks
on eight benchmarks \cite{DZT25}, Arp et al.\ measure the prevalence of ten
pitfalls across four venues \cite{AQPWPWCR22}, and Shafran et al.\ report that
model extraction has no formal definition of success at all \cite{SSEP24}. Balle,
Cherubin and Hayes define reconstruction against an informed adversary and quantify
over every reconstruction attack with no bound on it \cite{BCH22}, which is the
definitional form the critiques ask for, and it prices nobody. These works bound
nobody because their claims are about practice, or about what success means,
rather than about an adversary's cost.

\subsection{The auditor and the mechanism}

\AuditorBound{} works bound an auditor and \MechanismBound{} a mechanism, parties
that neither attack nor learn.

The auditor results are impossibilities about detection. Goldwasser, Kim,
Vaikuntanathan and Zamir plant a backdoor in a classifier so that no polynomial-time
distinguisher with oracle access can tell the backdoored model from an honestly
trained one \cite{GKVZ22}, and the bound sits on the distinguisher, so the
constrained party is whoever audits. Kalavasis et al.\ extend the guarantee to
white-box access under obfuscation, and note that the earlier work leaves that case
open \cite{KKOSVZ24}. Ball et al.\ show there exist language models for which no
efficient prompt filter works, under a construction in which the model shares a
secret key with the attacker \cite{BGGKRR26}. The claim is existential and about
constructed models.

The mechanism results price the defence itself. Bun, Chen and Vadhan separate
computational from statistical differential privacy by exhibiting a polynomial-time
mechanism, so the bound sits on the mechanism \cite{BCV16}. Dinur et al.\ bound the
space a differentially private mechanism may use while remaining sound under
adaptive queries \cite{DSWZ23}, and Ratliff and Vadhan treat the mechanism's own
running time as a leakage channel \cite{RV24}. All three put the resource bound on
the system rather than on anyone attacking it.

\subsection{Where the two halves meet, and do not}

Fine-grained cryptography is where the relation between two parties' work has a
name and a parameter \cite{BC22,LLW19,BRSV17}, and Degwekar, Vaikuntanathan and
Vasudevan restrict both parties to a circuit class instead \cite{DVV16}. It contains
no machine learning. The proof-of-learning line states the same relation with the
exponent fixed at one, as a requirement its authors could not prove
\cite{JYCDTCP21}, a gap their follow-up confirms \cite{FJTYCDCP23}. The
\NumRelativeAttack{} attack works of Table~\ref{tab:relative} each state such a
relation, and none writes it in the fine-grained vocabulary or cites it.

The learning-theory exponents are the closest the machine learning side comes to a
priced adversary. The information exponent \cite{AGJ21}, the leap \cite{ABM23} and
the generative exponent \cite{DPLB24,DLB25} each give a sample threshold defined
against an algorithm class.

Two boundaries on those thresholds are quoted more often than they now hold. The
thresholds were stated for detection rather than recovery, as the single-index
result says of itself \cite{DPLB24}. For polynomial links Lee, Oko, Suzuki and Wu
carry weak recovery through to strong recovery at
$\tilde\Theta(d^{(p^*-1) \vee 1})$ samples, governed by the generative exponent
\cite{LOSW24}, and Glasgow reaches $o(1)$ population error on the Boolean parity
instance in $\tilde\Theta(d)$ \cite{Gla24}. And they were stated for single-pass
training, which batch reuse was reported to escape \cite{DTAPZK24}. Reuse instead
attains the generative exponent in place of the information exponent, because two
steps on one minibatch pass the label through the activation's derivative
\cite{LOSW24,TME25}, and for the plain quadratic activation full-batch reuse gains
nothing at all \cite{KJWSM26}. A security claim resting on either boundary has to
say which version it means.

\subsection{Training compute as a priced resource}

The corpus holds \NumScaling{} works on neural scaling laws, and \NumFlops{} of
them are priced in a budget of floating-point operations spent on training, the
rest in the sample or parameter count that a fixed budget buys. Kaplan et al.\ fit
loss against model size, data and compute \cite{KMHBCCGRWA20}, and Hoffmann et al.\
correct the allocation between model size and data at a fixed budget
\cite{HBMBCRCHWCHNMDDGOSERVS22}. Ho et al.\ measure how much of the observed
improvement comes from better algorithms rather than more compute
\cite{HBEORGATS24}.

None of these works concerns an adversary, and all of them bound the trainer. They
matter to this survey because they establish that the unit exists and is measured.
Shavit's compute-monitoring proposal is the one work in the corpus that carries a
training-compute budget toward a security claim, and it states a hardness
requirement without a construction \cite{Sha23}. Its monitored quantity is
chip-hours rather than operations, which is a further step away from the unit the
scaling literature measures. The gap between a literature that measures training
compute precisely and a security literature that never states an adversary's budget
in it is the widest single gap this corpus contains.

One work crosses in the other direction. Haghifam, Smith and Ullman show a setting
where a membership inference attacker needs many more samples than the training
algorithm used \cite{HSU25}. The separation vanishes when the attacker knows the
covariance, so what it prices is ignorance of the data distribution rather than
attack difficulty as such.
